\chapter{Glossary of Important Concepts}
\markboth{Glossary of Important Concepts}{Glossary of Important Concepts}

These definitions distinguish geometric objects, physical models, and the guarantees obtained from a calculation. A definition does not establish that a particular golf model satisfies its assumptions.

\begin{description}[style=nextline,leftmargin=0pt,font=\bfseries,itemsep=0.5em]
\item[Articulated Body Algorithm]
A recursive forward-dynamics algorithm for a rigid-body tree. With bounded joint dimension, its arithmetic cost scales linearly with the number of bodies. Closed-loop constraints and contact require additional treatment; the tree algorithm alone does not solve every constrained contact problem.

\item[Configuration Space]
The set of configurations allowed by a model, including positions, orientations, and declared geometric constraints. It is a smooth manifold when suitable regularity conditions hold; redundant coordinates and constraint singularities require care. Velocities are additional state information.

\item[Contraction Analysis]
Analysis of convergence between trajectories under stated common inputs or closed-loop dynamics, in a specified metric and domain. A differential contraction inequality needs regularity, metric bounds, and domain/existence conditions to imply the corresponding finite-distance result. Contraction does not by itself prove input feasibility or a desired contact outcome.

\item[Control Contraction Metric (CCM)]
A positive-definite differential metric satisfying control-dependent contraction conditions that support feedback construction for feasible reference trajectories. The result depends on the theorem's domain, regularity, metric bounds, and controller construction. Actuator saturation and state constraints require explicit treatment.

\item[Degrees of Freedom (DOF)]
The number of locally independent coordinates needed to specify a configuration. Independent constraints can reduce it; redundant parameters can exceed it. The number of actuators and the dimension of the full state space are different quantities.

\item[Drift]
The evolution of the declared effective plant when the declared control channel is zero. It includes the configuration-velocity kinematics and all retained uncontrolled dynamics, such as gravity, inertial coupling, passive elasticity, and damping. Its meaning depends on the chosen state, control channel, and environment; zero control does not mean zero acceleration or zero velocity.

\item[Drift-Control Ratio (DCR)]
A model-based comparison of drift and bounded control authority in a declared acceleration or task-projected space:
\begin{equation*}
\operatorname{DCR}_{W,\mathcal U}(x)=
\frac{\|W a_{\mathrm{drift}}(x)\|_2}
{\sup_{u\in\mathcal U(x)}\|W B_a(x)u\|_2+\varepsilon} .
\end{equation*}
State the projection/scaling $W$, admissible input set $\mathcal U$, acceleration input map $B_a$, and regularizer $\varepsilon$, with compatible units. This scalar magnitude ratio omits directional alignment and does not by itself establish controllability, cancellation feasibility, or a universal coaching threshold.

\item[Euler Angles]
Three parameters describing an orientation using a declared sequence of rotations. Different sequences have different singular sets. Angle rates are related to angular velocity by a configuration-dependent map and are not generally its Cartesian components.

\item[Exponential Map]
For a Lie group, the map from a Lie-algebra element to the group element reached at parameter one along its generated one-parameter subgroup; for a matrix Lie group it is the matrix exponential. It need not be globally one-to-one. It is not general integration of an arbitrary time-varying velocity. The Riemannian exponential defined by geodesics is a related but distinct construction requiring a connection or metric.

\item[Flow]
The solution map of an ordinary differential equation where solutions exist uniquely. For an autonomous smooth vector field it forms a local one-parameter family with the composition property. Time-dependent systems use a two-time solution map; hybrid resets need not belong to a single smooth flow.

\item[Gimbal Lock]
A singularity of an Euler-angle rate-to-angular-velocity map. The parameterization loses local rank; the rigid body's physical rotational freedom does not disappear. Quaternions or overlapping coordinate charts handle orientation differently and have their own representation conventions.

\item[Jacobian]
The derivative of a map expressed in coordinates. For a task map $y=h(q)$, $\dot y=J(q)\dot q$. Acceleration also contains $\dot J\dot q$. Spatial and body Jacobians use different frame conventions, which must match the represented velocities and forces.

\item[Kinematics]
Description of position, orientation, velocity, and acceleration relationships without using forces to determine the motion. A kinematically allowed path is not automatically dynamically feasible under available actuation.

\item[Lagrangian Mechanics]
A formulation using a Lagrangian, often $L=T-U$, together with applied generalized forces and any constraints. For unconstrained coordinates, $d(\partial L/\partial\dot q)/dt-\partial L/\partial q=Q$. Dissipation, actuation, nonconservative loads, and contact must be specified; kinetic and potential energy alone do not determine a forced system's motion.

\item[Lie Algebra]
A vector space equipped with a bilinear antisymmetric bracket satisfying the Jacobi identity. For a Lie group, it is the tangent space at the identity with the induced bracket. Brackets of smooth vector fields are another important example; accessibility arguments require the appropriate system and theorem assumptions.

\item[Lie Group]
A smooth manifold that is also a group, with smooth multiplication and inversion. Examples include the rotation group $SO(3)$ and rigid-transformation group $SE(3)$. A choice of coordinates represents these objects without changing the group itself.

\item[Lyapunov Stability]
For an equilibrium, the property that sufficiently small initial perturbations remain small for all future times of the stated system. It does not require return to equilibrium. Asymptotic stability additionally requires attraction; orbital and finite-horizon stability statements refer to different targets and must be defined separately.

\item[Manifold]
A space locally described by Euclidean coordinates, with the usual topological regularity conditions. A smooth manifold also has smoothly compatible charts. Its dimension is local and does not depend on how many coordinates are used to embed or redundantly represent it.

\item[Screw Axis]
For a rigid-body twist with nonzero angular velocity, an instantaneous axis and pitch describing rotation together with translation along that axis. Pure translation is a limiting case, often represented as an axis at infinity; a finite unique rotational axis is then unavailable.

\item[State Space]
The set of variables sufficient to determine future evolution under specified inputs and a model. Mechanical models often include configuration and velocity or momentum, and may additionally include actuator states, flexible modes, estimators, or discrete contact modes.

\item[Tangent Bundle]
The union of tangent spaces $T_qQ$ over configurations $q\in Q$, with its bundle and smooth-manifold structure. For an $n$-dimensional smooth configuration manifold, $TQ$ has dimension $2n$. It represents configuration-velocity states for a basic mechanical model; momentum states instead lie in the cotangent bundle $T^*Q$.

\item[Tangent Space]
The vector space of instantaneous derivatives of smooth curves through a point on a manifold. Its vectors are local velocities, not finite displacements between arbitrary configurations. Applied differential constraints can restrict the admissible velocities further.

\item[Twist]
A six-component representation of an instantaneous rigid-body velocity field, combining angular and linear terms with a declared reference point and frame. In one frame, $v(r)=v_0+\omega\times(r-r_0)$. Spatial and body representations are related by the appropriate adjoint transformation; the linear term must not be confused with the velocity of an unspecified point.

\item[Wrench]
The dual force representation consisting of a moment about a declared reference point and a resultant force in a declared frame. With matching twist conventions, mechanical power is $\tau\cdot\omega+F\cdot v_0$. Changing point or frame requires the corresponding moment shift and dual transformation so that power is unchanged.

\item[Zero-Torque Counterfactual (ZTCF) Family]
Model interventions setting the declared applied generalized-control channel to zero while preserving the declared effective plant. Pointwise samples, stitched traces, forward trajectories, and branched trajectories answer different questions. A stitched pointwise acceleration trace is not automatically a feasible unforced trajectory.

\item[Zero-Velocity Counterfactual (ZVCF)]
The instantaneous acceleration after the declared velocity and control intervention:
\begin{equation*}
a_{\mathrm{ZVCF}}(q,z)=-M(q)^{-1}h(q,0,z_0;p) .
\end{equation*}
Specify the auxiliary-state reset $z_0$, retained loads, parameters, and contact mode. This separates the retained configuration-dependent loads under that intervention; it is not necessarily gravity alone, and it is not a forward simulation of a stationary body.
\end{description}
